\chapter{Conclusion and Future Work}
\label{chap:conclusion}
\label{chap:limitations}

% ===========================================================================
\section{Conclusion}
\label{sec:conclusion}

{\small The first goal, portability, was addressed by realising the framework as a single ROS~2 package~\cite{macenski2022ros2} with three nodes, namely the preprocessor, the Phase~1 semantic coordinator and the scene graph fuser, which exchange data with Hydra~\cite{hughes2022hydra} over standard topics (Sections~\ref{subsec:impl_deployment} and~\ref{subsec:impl_phase1}). The same package was run on the Jetson AGX Orin and on the Jetson AGX Thor (Section~\ref{sec:results_runtime}). Models and datasets are exchanged through configuration files only (Section~\ref{subsec:impl_portability}). Although the data recorded with the intended platform, the Unitree Go1 quadruped robot, proved unsuitable for the evaluation (Section~\ref{subsec:impl_go1_limitations}), the framework was run on real-world data from the office sequences of OpenLORIS-Scene~\cite{shi2020openloris}, which were added in this way without changing the code used for uHumans2~\cite{rosinol2021kimera}. These sequences were recorded on a wheeled robot with an Intel RealSense camera of the D400 series, the same sensor family as the D455 intended for the Go1, together with wheel odometry and a calibrated base-to-camera transform, and thus provide a comparable setting on real sensor data (Section~\ref{subsec:impl_datasets}).\par}

{\small The second goal combined spatial grounding, consistency on revisits and open-vocabulary identification. Class-agnostic object proposals from NanoSAM (Section~\ref{sec:results_segmentation}) are kept as persistent tracks whose slot IDs ground them as object nodes in Hydra's scene graph, so that semantic information and spatial map share one reference frame (Section~\ref{sec:method_tracking}). The evidence-based association (Section~\ref{subsec:method_association}) was deliberately made conservative: following the asymmetric-error principle, fragmenting an object into several tracks is preferred to merging two objects into one, since a merged track gives both objects a wrong geometry and label that cannot be repaired downstream, whereas fragmentation only duplicates an object (Sections~\ref{sec:method_tracking} and~\ref{sec:results_tracking}). Together with session persistence (Section~\ref{subsec:method_sessions}), this allows objects to be re-identified when an environment is revisited instead of being mapped again (Section~\ref{sec:results_tracking}). Objects are labelled by a vision-language model without a fixed list of classes (Section~\ref{sec:method_labelling}), and the evaluation showed that a compact model guided by a structured prompt is sufficient for this task, with the design of the prompt mattering more than the size of the model (Section~\ref{sec:results_vlm}).\par}

{\small For the third goal, each object carries separate information on its identity, its mobility and its continued presence (Section~\ref{sec:method_uncertainty}). Identity and mobility are determined once per track, when the object is labelled, and are represented by the confidences reported with them; objects without a usable label remain explicitly unclassified. The confidence in the continued presence of an object decays from the time of its last observation with a half-life that depends on its mobility class (Equation~\ref{eq:presence_decay}). As stated in Section~\ref{sec:method_uncertainty}, these quantities are confidence values and decay heuristics rather than calibrated probabilities, and they should be interpreted as relative measures of reliability.\par}

{\small For the fourth goal, the processing was divided into a fast per-frame path for segmentation, tracking and the input to Hydra, and an asynchronous path for semantic labelling, so that the computationally demanding vision-language model does not hold back the spatial map (Sections~\ref{subsec:impl_phase1} and~\ref{sec:results_runtime}). On the Jetson AGX Thor, geometry and tracking keep pace with the processed frames, segmentation remains the dominant cost of a frame, and moving from the Orin to the Thor roughly doubled the throughput without changes to the framework (Section~\ref{sec:results_runtime}). Semantic labelling does not keep pace on a first visit, when new objects arrive faster than the vision-language model can label them, but it catches up on revisits, since known objects keep their labels and hardly need further labelling calls (Sections~\ref{sec:results_runtime} and~\ref{sec:results_tracking}). Near-real-time operation was thus achieved for the spatial part of the graph, and for its semantic part once an environment is known.\par}

{\small Finally, the complete pipeline was run on two simulated environments, the office and apartment scenes of uHumans2~\cite{rosinol2021kimera}, and on two real-world sequences of OpenLORIS-Scene~\cite{shi2020openloris}, and produced a complete scene graph in each case (Sections~\ref{sec:results_tracking} to~\ref{sec:results_visual}). In all of them, Hydra organised the reconstructed environment into places and rooms, and the fused graph placed the tracked objects in this hierarchy with their labels, mobility classes, visual-state descriptions and relations (Section~\ref{sec:results_visual}). Moving between the four environments required only a change of configuration (Section~\ref{subsec:impl_portability}), which shows that the framework is compatible with both simulated and real-world data, as set out in the goals of this thesis.\par}

{\small The software architecture is designed so that the pipeline can benefit from future models without being rebuilt. Each model is accessed only through a small interface, one each for segmentation, retrieval and the vision-language model, and the implementation in use is selected in the configuration file (Section~\ref{subsec:impl_portability}). A vision-language model served through a standard chat-completions endpoint is added as a new model profile without any change to the code, and a segmentation model with a different interface requires only one new class, since all later stages consume the same mask format. This was confirmed during the evaluation: the newer Qwen3.5 model replaced the vision-language model of the baseline by changing only the model profile and the prompt, and improved the labelling without changes to the other components (Section~\ref{sec:results_vlm}). Since small vision-language and segmentation models continue to improve rapidly (Section~\ref{subsec:qwen}), the framework can adopt such models as they become available, with only a re-validation of their accuracy on the existing prompt.\par}

{\small Taken together, this thesis provides a framework for downstream risk assessment by building a spatially consistent scene graph, in which each tracked object appears as a persistent node with an open-vocabulary label, a mobility class, the confidences of both, a presence confidence and a short description of its visual state (Section~\ref{subsec:method_vlm_contract}). The objects are linked by contact and identity relations and assigned to places and rooms where the reconstructed geometry supports it (Section~\ref{subsec:method_graph_fusion}). The thesis therefore produces the risk-relevant information on which a dynamic risk assessment can operate; it does not itself compute a risk score.\par}


% ===========================================================================
\section{Limitations and Future Work}
\label{sec:limitations}
\label{sec:future_work}

{\small The framework depends heavily on the accuracy of the depth data, since depth is used to back-project the masks, to associate observations with tracks and to build Hydra's reconstruction (Section~\ref{subsec:impl_go1_limitations}). With active stereo cameras, depth quality depends strongly on lighting, reflective surfaces and the distance to the object, as the recordings with the Go1 showed (Section~\ref{subsec:impl_go1_limitations}). The reliability of the geometry could be improved by combining the camera with further sensors whose weaknesses differ from those of stereo depth. A LiDAR sensor is the most direct complement, since it measures distance actively by time of flight and is therefore largely independent of lighting and surface texture, while the camera continues to provide the colour information required for segmentation and labelling. Such a combination of complementary sensors is already common in automated driving, where cameras, LiDAR and radar are used together. For the present framework, LiDAR points projected into the camera image could supply or verify the depth of each mask, and LiDAR-based odometry could complement the wheel odometry used so far, so that both the object geometry and the map would rely less on a single sensor.\par}

{\small Segmentation limits both the throughput of the pipeline, that is, the number of frames it can process per second, and the quality of the scene graph. Every processed frame is segmented on the per-frame path, and segmentation accounts for the largest share of the processing time of a frame (Section~\ref{sec:results_runtime}). At the same time, all subsequent steps, from the 3D observations to the crops passed to the vision-language model, are derived from its masks (Sections~\ref{sec:method_tracking} and~\ref{sec:method_labelling}): objects that are missed do not enter the graph, a mask that covers two touching surfaces cannot be separated by the tracker, and an object that is segmented inconsistently across views can be represented by several tracks and thus by several nodes (Sections~\ref{sec:results_segmentation} and~\ref{sec:results_tracking}). Frames could be filtered before they are passed to the segmentation model, so that only frames that add new information, for example after a sufficient change of the viewpoint, are segmented. Newer segmentation models should be chosen for accuracy as well as speed, since models that separate instances more reliably and produce consistent masks across views would reduce merged masks and fragmentation; the step from SAM with a ViT-B encoder to NanoSAM already showed that both can improve together (Section~\ref{sec:results_segmentation}). The decisive step, however, would be a system that can follow objects that have already been identified from frame to frame without segmenting every frame anew, so that the segmentation model is needed only to discover new objects.\par}

{\small Each object is represented by an axis-aligned bounding box in the map frame, and the association compares these boxes to decide whether an observation belongs to a track (Sections~\ref{subsec:method_proposals} and~\ref{subsec:method_association}). Such boxes are cheap to compute and were sufficient for the evaluated scenes, but they describe the shape of an object only coarsely. For an elongated object that lies diagonally to the axes, the box encloses much more space than the object itself. Neighbouring objects may then appear to overlap with it, which makes the association less reliable. Future work could address this in two ways. The first is to use oriented bounding boxes, which are rotated to follow the main directions of the object. They fit the object closely at little extra cost, but their orientation must be kept stable across observations, which is difficult for nearly symmetric objects. The second is to track objects by their reconstructed surface, for example by comparing new observations with the part of Hydra's mesh that belongs to the object. This describes the object itself rather than an enclosing box, but requires more computation. Both approaches share a more basic limitation, however: the camera sees an object from only one side at a time, so its far side remains unknown. A large object that is first seen from one side and later from the opposite side may therefore not be recognised as the same object and is then represented by two tracks.\par}

{\small In the labelling stage, retrieval was disabled in the evaluated configuration, since the persistent tracks already preserve the identity of known objects (Sections~\ref{sec:results_vlm} and~\ref{sec:results_visual}), and a single call to the vision-language model provides both the label and the risk-relevant visual-state description (Section~\ref{subsec:method_vlm_contract}). An alternative would be to separate the two into two calls and to enable retrieval for labelling. Retrieval would then resolve part of the labels and reduce the load on the vision-language model, at the cost of additional retrieval load, whereas the description would still be requested for every object, since it depends on the state and context of the individual instance rather than on its label alone. Which of the two configurations achieves the higher throughput and the more reliable labels can only be decided by a comparative study, which should precede any change to the current design.\par}

{\small In the fused scene graph, tracks are removed only for dynamic objects; static objects and objects of unknown mobility are never withdrawn and only lose presence confidence over time (Section~\ref{subsec:method_static_dynamic}). The framework therefore does not cover the case in which a static object is moved while the camera is not observing it. When the area is revisited, the object remains in the scene graph at its old position, because its absence from a location that is in view is not used as evidence that it has gone. A mechanism that also uses such negative evidence could correct this, for example a score for each track that increases with repeated observations and decreases whenever the position of the object is in view but the object is not observed, so that the track is removed once the score falls below a threshold.\par}

{\small The confidence values attached to each object are not calibrated probabilities. The label and mobility confidences are reported by the vision-language model, and the presence confidence follows a decay heuristic with fixed half-lives, so that all of them can only be interpreted as relative measures of reliability (Section~\ref{sec:method_uncertainty}). Future work could calibrate these values against observed outcomes, for example by relating the reported label confidences to the measured labelling accuracy, and by deriving the half-lives of the presence decay from how often objects of each mobility class actually change their position, so that downstream risk assessment can treat them as probabilities.\par}

{\small Finally, the throughput is bounded by the embedded platform, on which segmentation, Hydra and the vision-language model share one GPU (Section~\ref{sec:results_runtime}). The move from the Jetson AGX Orin to the Jetson AGX Thor already raised the throughput without changes to the framework, and more capable embedded hardware can be expected to raise it further. On a given platform, the throughput could also be improved by optimising the allocation of resources such as memory and processor cores among the components.\par}

{\small None of these limitations calls the overall design into question. Most of them concern individual components, such as the depth sensor, the segmentation model or the representation of object extents, which can be improved or exchanged without changing the structure of the pipeline. The framework developed in this thesis therefore provides a foundation on which more capable sensors, models and hardware can be built, bringing a continuously updated, spatially consistent description of the environment closer to practical use on mobile robots.\par}